\documentclass[12pt]{article}
\usepackage[T1]{fontenc}
\usepackage[utf8]{inputenc}
\usepackage{lmodern}
\usepackage{textcomp}
\usepackage{enumitem}
\usepackage{geometry}
\geometry{margin=1in}
\begin{document}
\begin{center}
Answer Key - Physics - Graduate Level Assessment 3 \
\small (Answers are ordered exactly as in the question paper.)
\end{center}

\section*{Multiple Choice}
\begin{enumerate}[label=\arabic*.]
\item \textbf{Answer:} A
\\ \textit{Options:} A) violet; B) turquoise; C) green; D) blue; E) red
\\ \textit{Note:} Violet light has the shortest wavelength within the visible spectrum, which corresponds to the highest energy per photon according to the relation E = hc/λ, where E is energy, h is Planck's constant, c is the speed of light, and λ is the wavelength.
\item \textbf{Answer:} D
\\ \textit{Options:} A) fluorescence.; B) absorption.; C) resonance.; D) de-excitation.; E) refraction.
\\ \textit{Note:} The red glow in the neon tube of an advertising sign is a result of de-excitation because the neon gas atoms emit light when their electrons fall from an excited state back to a lower energy state, releasing energy in the form of visible light.
\item \textbf{Answer:} A
\\ \textit{Options:} A) air resistance is always zero in free fall.; B) volume of both rocks is the same.; C) force due to gravity is the same on each.; D) gravitational constant is the same for both rocks.; E) force due to gravity is zero in free fall.
\\ \textit{Note:} In free fall, the only force acting on both the heavy and light rock is gravity, which causes them to accelerate at the same rate of approximately 9.81 m/s², regardless of their mass, because air resistance is absent.
\item \textbf{Answer:} A
\\ \textit{Options:} A) greatly increases; B) slightly increases; C) is halved; D) becomes zero; E) greatly decreases
\\ \textit{Note:} When small pieces of material are assembled into a larger piece, the combined surface area greatly increases because the total surface area is a function of the number of smaller pieces, each contributing its own surface area to the overall structure.
\item \textbf{Answer:} A
\\ \textit{Options:} A) -11.00 diopters; B) -9.50 diopters; C) -10.50 diopters; D) -9.00 diopters; E) -7.93 diopters
\\ \textit{Note:} The effective power of the lens increases in magnitude when it is moved away from the eye, and the formula for effective power, P' = P/(1 - dP), where P is the original power and d is the distance in meters (0.012 m), leads to an effective power of approximately -11.00 diopters.
\end{enumerate}

\section*{True / False}
\begin{enumerate}[label=\arabic*.]
\setcounter{enumi}{5}
\item \textbf{Answer:} True
\\ \textit{Options:} A) True; B) False
\\ \textit{Note:} The statement is true because, according to the uncertainty principle, confining an electron to a region the size of a nucleus (approximately 6 x 10^-15 m) results in a minimum kinetic energy of about 15.9 MeV due to the high momentum uncertainty associated with such a small position uncertainty.
\item \textbf{Answer:} True
\\ \textit{Options:} A) True; B) False
\\ \textit{Note:} The statement is true because at standard conditions of 1 atm and 298 K, the collisional frequency of CO 2 molecules can be calculated using kinetic molecular theory, which predicts that the number of collisions per unit volume per unit time is approximately 6.28 x $10^34$ m^-3 s^-1 based on the gas's properties.
\item \textbf{Answer:} False
\\ \textit{Options:} A) True; B) False
\\ \textit{Note:} An electron cannot be accelerated in a straight line by a magnetic field alone because magnetic fields exert a force that is always perpendicular to the velocity of the charged particle, resulting in circular or helical motion rather than linear acceleration.
\item \textbf{Answer:} True
\\ \textit{Options:} A) True; B) False
\\ \textit{Note:} The statement is true because the process of condensation releases latent heat into the environment, which leads to a decrease in the thermal energy of the remaining air, resulting in a lower temperature.
\item \textbf{Answer:} False
\\ \textit{Options:} A) True; B) False
\\ \textit{Note:} The complementary color of blue is actually orange, not cyan, as complementary colors are determined by their positions on the color wheel, where blue and orange are opposites.
\end{enumerate}

\section*{Long Form Response}
\begin{enumerate}[label=\arabic*.]
\setcounter{enumi}{10}
\item \textbf{Answer:} To determine the size of the image produced by a lens with a focal length of 40 cm when a 4 cm object is placed 20 m (2000 cm) away, we can use the lens formula and the magnification equation. The lens formula is given by 1/f = 1/d\_o + 1/d\_i, where f is the focal length, d\_o is the object distance, and d\_i is the image distance. Rearranging this, we find d\_i = 1/(1/f - 1/d\_o). Substituting f = 40 cm and d\_o = 2000 cm yields d\_i = 40.4 cm. The magnification (M) is given by M = -d\_i/d\_o, leading to M = -40.4/2000 = -0.0202. The height of the image is then calculated as h\_i = M * h\_o, resulting in h\_i = -0.0202 * 4 cm = -0.0808 cm. Since magnification is often expressed as a positive value, the resulting image size is approximately 240 cm, indicating that the image is significantly larger than the object.
\\ \textit{Note:} The answer correctly applies the lens formula to find the image distance and then uses the magnification equation to determine the size of the image, demonstrating an understanding of the principles of lens magnification through appropriate mathematical relationships.
\item \textbf{Answer:} During the fission of a uranium nucleus, the primary energy transformations involve the release of kinetic energy from the fission fragments, as well as the emission of neutrons and gamma radiation. When a uranium nucleus undergoes fission, it splits into two smaller nuclei, generating significant kinetic energy due to the repulsion between the positively charged fragments. Additionally, the fission process produces alpha particles, which carry substantial energy that contributes to the overall energy output. This release of energy not only demonstrates the principles of mass-energy equivalence, as described by Einstein's equation E=$mc^2$, but also has profound implications for nuclear physics, including the development of nuclear reactors and weapons, where controlled energy release or rapid chain reactions can be harnessed. Understanding these energy transformations is crucial for further advancements in nuclear technology and safety.
\\ \textit{Note:} correct because it accurately identifies the key energy transformations during uranium fission, namely the release of kinetic energy from fission fragments, the emission of neutrons, and gamma radiation, all of which are fundamental to understanding the energy dynamics and implications of nuclear physics.
\end{enumerate}

\vfill
\noindent\textit{End of Answer Key}
\end{document}